% !PRINTED 34
\weberpage{34}

Wenn eine Function durch eine andere theilbar ist, so ist der Grad des Quotienten gleich der Differenz des Grades des Dividenden und des Grades des Divisors. Ersterer kann also nicht kleiner sein als letzterer. Sind die Grade gleich, so ist der Quotient eine Constante, und daraus folgt:

\gesperrt{7. Wenn von zwei ganzen rationalen Functionen gleichen Grades die eine durch die andere theilbar ist, so unterscheiden sie sich nur durch einen constanten Factor von einander, und es ist auch die zweite durch die erste theilbar.}

\gesperrt{8. Nach \S~4 ist die nothwendige und hinreichende Bedingung dafür, dass die Function $f(x)$ durch die lineare Function $x-\alpha$ theilbar ist, die dass $f(\alpha)=0$ sei.}

\begin{center}
{\large \S.\ 6.}

\vspace{0.5em}
{\large\gesperrt{Grösster gemeinschaftlicher Theiler.}}
\end{center}

Es ist eine Aufgabe von fundamentaler Bedeutung, zu entscheiden, ob zwei ganze rationale Functionen ausser den Constanten noch einen anderen gemeinschaftlichen Theiler haben. Man findet die Lösung durch den \gesperrt{Algorithmus des grössten gemeinschaftlichen Theilers} ganz in derselben Weise, wie die entsprechende Frage in Bezug auf die Theilbarkeit ganzer Zahlen beantwortet wird. (S. die Einleitung.)

Es seien
\begin{equation}
 f(x)=A,\qquad \varphi(x)=A'
\tag{1}
\end{equation}
zwei gegebene Functionen; der Grad von $\varphi(x)$ möge niedriger oder wenigstens nicht höher als der von $f(x)$ sein.

Wir dividiren $A$ durch $A'$ und bezeichnen den Rest, dessen Grad niedriger ist als der Grad von $A'$, mit $A''$, also:
\[
 A=Q'A'+A'';
\]
nun dividiren wir $A'$ durch $A''$ und bezeichnen den Rest mit $A'''$, und fahren so fort in der Bildung der Functionenreihe
\begin{equation}
 A,A',A'',A''',\ldots,
\tag{2}
\end{equation}
deren Grade $n,n',n'',n'''\ldots$ immer abnehmen und folglich nach einer endlichen Anzahl von Divisionen auf Null heruntergehen. Der letzte Rest vom Grade Null, der also eine Constante ist, sei $A^{(\nu)}$. Dann haben wir die Kette der Gleichungen:
%CONT>>
